\newpage
\thispagestyle{empty}

\begin{center}
    \Large\textbf{摘要}
\end{center}

\vspace{1cm}

神经网络彻底改变了人工智能领域，使计算机视觉、自然语言处理、语音识别等众多领域实现了重大突破。本教材全面介绍了神经网络，从基本的数学概念到最新的网络架构和训练技术。

本书共分十三章，内容涵盖：

\begin{enumerate}
    \item \textbf{数学基础：} 线性代数（SVD、特征分解、条件数）、微积分（泰勒展开、海森矩阵、凸性）、概率论（最大似然估计、最大后验估计、信息论）和优化理论。

    \item \textbf{激活函数：} 对 Sigmoid、tanh、ReLU 及其变体（Leaky ReLU、PReLU、ELU）、Swish/SiLU、Mish、GELU、StarReLU 进行了全面分析，并给出了通用近似定理的证明和选择标准指导。

    \item \textbf{损失函数：} 详细探讨了回归损失（MSE、Huber、MAE）、分类损失（交叉熵、Focal Loss、标签平滑）以及对比损失和 Dice 损失等任务特定目标函数。

    \item \textbf{神经网络架构：} 深度神经网络（DNN）、卷积神经网络（CNN）和循环神经网络（RNN/LSTM/GRU），包含数学推导和实际示例。

    \item \textbf{核心操作：} 卷积（标准卷积、膨胀卷积、深度可分离卷积、转置卷积）、池化和归一化操作（Batch Norm、Layer Norm、Group Norm、Instance Norm），具有数学严谨性。

    \item \textbf{高级机制：} 跳跃连接（残差连接、密集连接）、通道注意力（SE 模块、CBAM）和空间注意力机制，包括梯度流的理论分析。

    \item \textbf{优化：} 梯度下降变体（SGD、动量、Nesterov）、自适应优化器（AdaGrad、RMSProp、Adam、AdamW、AMSGrad）以及学习率调度器（余弦退火、周期性调度、热重启）。

    \item \textbf{经典架构：} 按时间顺序介绍从 LeNet-5（1998）和 AlexNet（2012）到 VGGNet、GoogLeNet/Inception、ResNet、DenseNet、MobileNet、SENet、EfficientNet、Vision Transformer（ViT，2020）、U-Net 和 ConvNeXt。

    \item \textbf{计算分析：} 全连接层、卷积层、注意力层和归一化层的参数量统计和 FLOPs 计算，以及模型压缩和缩放定律。

    \item \textbf{训练流程与框架：} 实际训练流程、正则化技术（Dropout、权重衰减、早停）、超参数调优、混合精度训练，以及 PyTorch 实现与 CUDA 概念。

    \item \textbf{数据集处理：} 数据收集、标注、增强策略（Mixup、CutMix、AutoAugment、RandAugment）、划分策略，以及图像、文本和音频的数据加载管道。

    \item \textbf{Transformer 架构：} 自注意力机制、多头注意力、位置编码（正弦编码、RoPE、ALiBi）、编码器-解码器架构、BERT、GPT，以及现代大语言模型（包括 LLaMA 和缩放定律）的概述。

    \item \textbf{现代训练技术：} 预训练和微调范式、参数高效微调（LoRA、适配器）、对比学习（SimCLR、InfoNCE）、掩码自编码（MAE），以及灾难性遗忘和多任务学习的进阶主题。
\end{enumerate}

每章都包含详细的数学推导、直观的解释、PyTorch 代码示例、自我评估练习以及经典和最新研究论文的参考文献。完成本教材后，读者将能够为各种实际应用设计、实现、训练和分析神经网络。

\vspace{2cm}

\textbf{关键词：} 神经网络、深度学习、机器学习、CNN、RNN、LSTM、Transformer、注意力机制、自注意力、优化、反向传播、PyTorch、计算机视觉、自然语言处理、大语言模型、迁移学习、微调

\newpage

